\chapter{Modelo de Cronograma de TCC}

O cronograma abaixo é apenas um modelo. Cada curso e cada estudante devem adaptá-lo ao calendário acadêmico, à modalidade de TCC e às exigências de orientação.

\begin{longtable}{p{0.36\linewidth}p{0.12\linewidth}p{0.12\linewidth}p{0.12\linewidth}p{0.12\linewidth}}
\toprule
Atividade & Mês 1 & Mês 2 & Mês 3 & Mês 4 \\
\midrule
\endfirsthead
\toprule
Atividade & Mês 1 & Mês 2 & Mês 3 & Mês 4 \\
\midrule
\endhead
Definição do tema & & & & \\
Delimitação e problema & & & & \\
Objetivos e justificativa & & & & \\
Levantamento bibliográfico & & & & \\
Projeto ou plano de trabalho & & & & \\
Consulta ética e autorizações & & & & \\
Submissão ao sistema competente, quando aplicável & & & & \\
Coleta ou seleção de dados & & & & \\
Análise dos dados & & & & \\
Redação das seções textuais & & & & \\
Revisão com orientador & & & & \\
Normalização & & & & \\
Entrega para banca & & & & \\
Correções finais & & & & \\
\bottomrule
\end{longtable}

\section{Orientações de uso}

O cronograma deve ser construído de trás para frente a partir da data final de entrega. Reserve tempo para atrasos, revisão e ajustes. Se houver comitê de ética, autorização institucional ou coleta externa, amplie o prazo.

Não inicie recrutamento ou coleta antes das autorizações e aprovações aplicáveis. Como esses prazos podem variar, registre dependências e alternativas em vez de concentrar todas as etapas em um único mês.
